Je werkt als junior systeembeheerder. Een server houdt gebeurtenissen bij in een CSV-bestand. De servicedesk wil alleen bepaalde meldingen ontvangen.

Je gaat daarom een Python-programma maken dat:

1. een CSV-bestand inleest;
2. de regels controleert;
3. alleen meldingen die aan een bepaalde voorwaarde voldoen selecteert;
4. de geselecteerde gegevens omzet naar JSON;
5. het resultaat opslaat in een JSON-bestand.

### Inputbestand

Maak een bestand `serverlog.csv`:

```csv
datum,tijd,server,gebruikersnaam,type,status,bericht
2026-09-15,08:01,SRV01,jansen,login,OK,Succesvol ingelogd
2026-09-15,08:03,SRV01,devries,login,FAIL,Verkeerd wachtwoord
2026-09-15,08:05,SRV02,admin,service,FAIL,Webserver gestopt
2026-09-15,08:07,SRV01,pietersen,login,OK,Succesvol ingelogd
2026-09-15,08:09,SRV03,admin,disk,WARNING,Schijfruimte onder 20 procent
2026-09-15,08:11,SRV02,admin,service,OK,Webserver gestart
2026-09-15,08:13,SRV03,jansen,login,FAIL,Account geblokkeerd
```

### Opdracht

Schrijf een Python-programma `logfilter.py`.

Het programma moet het CSV-bestand lezen en **alleen de regels selecteren waarbij `status` gelijk is aan `FAIL`**.

De geselecteerde regels moeten vervolgens worden opgeslagen in:

```text
resultaat.json
```

De JSON moet bijvoorbeeld deze structuur hebben:

```json
[
    {
        "datum": "2026-09-15",
        "tijd": "08:03",
        "server": "SRV01",
        "gebruikersnaam": "devries",
        "type": "login",
        "status": "FAIL",
        "bericht": "Verkeerd wachtwoord"
    },
    {
        "datum": "2026-09-15",
        "tijd": "08:05",
        "server": "SRV02",
        "gebruikersnaam": "admin",
        "type": "service",
        "status": "FAIL",
        "bericht": "Webserver gestopt"
    }
]
```

### Eisen

Het programma moet:

* gebruikmaken van de Python-module `csv`;
* gebruikmaken van de Python-module `json`;
* het CSV-bestand met Python openen;
* de inhoud van iedere regel controleren;
* alleen `FAIL`-meldingen selecteren;
* de geselecteerde gegevens bewaren in een Python-lijst;
* de lijst omzetten naar JSON;
* `resultaat.json` automatisch aanmaken;
* JSON netjes leesbaar formatteren met bijvoorbeeld `indent=4`.

### Uitbreiding 1 — filter instellen

Maak het programma flexibeler.

Laat de gebruiker kiezen waarop gefilterd wordt:

```text
Waarop wil je filteren?
1. FAIL
2. WARNING
3. OK

Maak een keuze:
```

Het programma gebruikt vervolgens de keuze als filter.

### Uitbreiding 2 — filter op server

Laat de gebruiker ook een server opgeven:

```text
Welke server wil je bekijken?
SRV01
```

Het programma geeft dan alleen meldingen van die server terug.

Bijvoorbeeld:

```text
Server: SRV01
Status: FAIL
```

### Uitbreiding 3 — combinatie van filters

Laat de gebruiker zowel een server als status invoeren:

```text
Server: SRV01
Status: FAIL
```

Alleen regels die aan **beide voorwaarden** voldoen worden naar JSON geschreven.

### Uitbreiding 4 — statistiek

Laat het programma vóór het maken van het JSON-bestand melden:

```text
Aantal regels gelezen: 7
Aantal regels geselecteerd: 3
Resultaat opgeslagen in resultaat.json
```

## Denkvragen voor de student

Beantwoord na afloop:

1. Wat is het verschil tussen CSV en JSON?
2. Waarom is een CSV-bestand handig als invoerbestand?
3. Waarom gebruiken we een Python `list` voor de gefilterde gegevens?
4. Wat doet `csv.DictReader()`?
5. Wat doet `json.dump()`?
6. Wat is het voordeel van `indent=4`?
7. Wat gebeurt er als het CSV-bestand niet bestaat?
8. Hoe zou je het programma aanpassen zodat alleen `WARNING`-meldingen worden weergegeven?

## Eindopdracht

Maak uiteindelijk een programma waarbij de gebruiker zelf kan aangeven:

```text
=== Server Log Filter ===

CSV-bestand: serverlog.csv

Server (bijvoorbeeld SRV01 of ALL): SRV01
Status (OK, WARNING, FAIL of ALL): FAIL

Aantal regels gelezen: 7
Aantal regels gevonden: 1

JSON-bestand: resultaat.json
```

De docent kan vervolgens controleren of `resultaat.json` uitsluitend de juiste records bevat.

**Extra leerzaam voor MBO 4:** laat studenten eerst zelf een oplossing maken zonder voorbeeldcode. Daarna kunnen ze `DictReader`, `json.dump()` en foutafhandeling toevoegen. Zo blijft de opdracht vooral gericht op **probleemoplossend programmeren**, in plaats van alleen code overtypen.

